\section{Projective bundles}
\label{section-projective-bundle}

\noindent
Let $S$ be a scheme.
Let $\mathcal{E}$ be a quasi-coherent sheaf of $\mathcal{O}_S$-modules.
By Modules, Lemma \ref{modules-lemma-whole-tensor-algebra-permanence}
the symmetric algebra $\text{Sym}(\mathcal{E})$ of
$\mathcal{E}$ over $\mathcal{O}_S$
is a quasi-coherent sheaf of $\mathcal{O}_S$-algebras.
Note that it is generated in degree $1$ over $\mathcal{O}_S$.
Hence it makes sense to apply the construction of the
previous section to it, specifically Lemmas
\ref{lemma-relative-proj} and \ref{lemma-apply-relative}.

\begin{definition}
\label{definition-projective-bundle}
Let $S$ be a scheme. Let $\mathcal{E}$ be a quasi-coherent
$\mathcal{O}_S$-module\footnote{The reader may expect here
the condition that $\mathcal{E}$ is finite locally free. We do not
do so in order to be consistent with
\cite[II, Definition 4.1.1]{EGA}.}.
We denote
$$
\pi :
\mathbf{P}(\mathcal{E}) = \underline{\text{Proj}}_S(\text{Sym}(\mathcal{E}))
\longrightarrow
S
$$
and we call it the {\it projective bundle associated to $\mathcal{E}$}.
The symbol $\mathcal{O}_{\mathbf{P}(\mathcal{E})}(n)$
indicates the invertible $\mathcal{O}_{\mathbf{P}(\mathcal{E})}$-module
of Lemma \ref{lemma-apply-relative} and is called the $n$th
{\it twist of the structure sheaf}.
\end{definition}

\noindent
According to Lemma \ref{lemma-glue-relative-proj-twists} there are
canonical $\mathcal{O}_S$-module homomorphisms
$$
\text{Sym}^n(\mathcal{E})
\longrightarrow
\pi_*\mathcal{O}_{\mathbf{P}(\mathcal{E})}(n)
\quad\text{equivalently}\quad
\pi^*\text{Sym}^n(\mathcal{E})
\longrightarrow
\mathcal{O}_{\mathbf{P}(\mathcal{E})}(n)
$$
for all $n \geq 0$. In particular, for $n = 1$ we have
$$
\mathcal{E}
\longrightarrow
\pi_*\mathcal{O}_{\mathbf{P}(\mathcal{E})}(1)
\quad\text{equivalently}\quad
\pi^*\mathcal{E}
\longrightarrow
\mathcal{O}_{\mathbf{P}(\mathcal{E})}(1)
$$
and the map $\pi^*\mathcal{E} \to \mathcal{O}_{\mathbf{P}(\mathcal{E})}(1)$
is a surjection by Lemma \ref{lemma-apply-relative}.
This is a good way to remember how we have normalized
our construction of $\mathbf{P}(\mathcal{E})$.

\medskip\noindent
Warning: In some references the scheme $\mathbf{P}(\mathcal{E})$
is only defined for $\mathcal{E}$ finite locally free on $S$.
Moreover sometimes $\mathbf{P}(\mathcal{E})$ is actually defined as our
$\mathbf{P}(\mathcal{E}^\vee)$ where $\mathcal{E}^\vee$
is the dual of $\mathcal{E}$ (and this is done only when $\mathcal{E}$ is
finite locally free).

\medskip\noindent
Let $S$, $\mathcal{E}$, $\mathbf{P}(\mathcal{E}) \to S$ be as in
Definition \ref{definition-projective-bundle}. Let $f : T \to S$
be a scheme over $S$. Let $\psi : f^*\mathcal{E} \to \mathcal{L}$
be a surjection where $\mathcal{L}$ is an invertible $\mathcal{O}_T$-module.
The induced graded $\mathcal{O}_T$-algebra map
$$
f^*\text{Sym}(\mathcal{E}) = \text{Sym}(f^*\mathcal{E}) \to
\text{Sym}(\mathcal{L}) = \bigoplus\nolimits_{n \geq 0} \mathcal{L}^{\otimes n}
$$
corresponds to a morphism
$$
\varphi_{\mathcal{L}, \psi} : T \longrightarrow \mathbf{P}(\mathcal{E})
$$
over $S$ by our construction of the relative Proj as the scheme representing
the functor $F$ in Section \ref{section-relative-proj}. On the other hand,
given a morphism $\varphi : T \to \mathbf{P}(\mathcal{E})$ over $S$
we can set $\mathcal{L} = \varphi^*\mathcal{O}_{\mathbf{P}(\mathcal{E})}(1)$
and $\psi : f^*\mathcal{E} \to \mathcal{L}$ equal to the pullback
by $\varphi$ of the canonical surjection
$\pi^*\mathcal{E} \to \mathcal{O}_{\mathbf{P}(\mathcal{E})}(1)$.
By Lemma \ref{lemma-apply-relative} these constructions
are inverse bijections between the set of isomorphism classes of pairs
$(\mathcal{L}, \psi)$ and the set of morphisms
$\varphi : T \to \mathbf{P}(\mathcal{E})$ over $S$.
Thus we see that $\mathbf{P}(\mathcal{E})$ represents the functor
which associates to $f : T \to S$ the set of $\mathcal{O}_T$-module
quotients of $f^*\mathcal{E}$ which are locally free of rank $1$.

\begin{example}[Projective space of a vector space]
\label{example-projective-space}
Let $k$ be a field. Let $V$ be a $k$-vector space. The corresponding
{\it projective space} is the $k$-scheme
$$
\mathbf{P}(V) = \text{Proj}(\text{Sym}(V))
$$
where $\text{Sym}(V)$ is the symmetric algebra on $V$ over $k$.
Of course we have $\mathbf{P}(V) \cong \mathbf{P}^n_k$ if $\dim(V) = n + 1$
because then the symmetric algebra on $V$ is isomorphic to a polynomial
ring in $n + 1$ variables. If we
think of $V$ as a quasi-coherent module on $\Spec(k)$, then $\mathbf{P}(V)$
is the corresponding projective space bundle over $\Spec(k)$. By the
discussion above a $k$-valued point $p$ of $\mathbf{P}(V)$ corresponds to
a surjection of $k$-vector spaces $V \to L_p$ with $\dim(L_p) = 1$.
More generally, let $X$ be a scheme over $k$, let $\mathcal{L}$ be an
invertible $\mathcal{O}_X$-module, and let
$\psi : V \to \Gamma(X, \mathcal{L})$ be a $k$-linear map
such that $\mathcal{L}$ is generated as an $\mathcal{O}_X$-module
by the sections in the image of $\psi$. Then the discussion above
gives a canonical morphism
$$
\varphi_{\mathcal{L}, \psi} : X \longrightarrow \mathbf{P}(V)
$$
of schemes over $k$ such that there is an isomorphism
$\theta : \varphi_{\mathcal{L}, \psi}^*\mathcal{O}_{\mathbf{P}(V)}(1)
\to \mathcal{L}$ and such that $\psi$ agrees with the composition
$$
V \to
\Gamma(\mathbf{P}(V), \mathcal{O}_{\mathbf{P}(V)}(1))
\to
\Gamma(X, \varphi_{\mathcal{L}, \psi}^*\mathcal{O}_{\mathbf{P}(V)}(1))
\to
\Gamma(X, \mathcal{L})
$$
See Lemma \ref{lemma-invertible-map-into-proj}. If
$V \subset \Gamma(X, \mathcal{L})$ is a subspace, then we will
denote the morphism constructed above simply as
$\varphi_{\mathcal{L}, V}$.
If $\dim(V) = n + 1$ and we choose a basis $v_0, \ldots, v_n$ of $V$
then the diagram
$$
\xymatrix{
X \ar@{=}[d] \ar[rr]_{\varphi_{\mathcal{L}, \psi}} & &
\mathbf{P}(V) \ar[d]^{\cong} \\
X \ar[rr]^{\varphi_{(\mathcal{L}, (s_0, \ldots, s_n))}} & &
\mathbf{P}^n_k
}
$$
is commutative, where $s_i = \psi(v_i) \in \Gamma(X, \mathcal{L})$, where
$\varphi_{(\mathcal{L}, (s_0, \ldots, s_n))}$
is as in Section \ref{section-projective-space},
and where the right vertical arrow corresponds
to the isomorphism $k[T_0, \ldots, T_n] \to \text{Sym}(V)$ sending
$T_i$ to $v_i$.
\end{example}

\begin{example}
\label{example-projective-bundle}
The map $\text{Sym}^n(\mathcal{E}) \to
\pi_*(\mathcal{O}_{\mathbf{P}(\mathcal{E})}(n))$
is an isomorphism if $\mathcal{E}$ is locally free, but in general
need not be an isomorphism. In fact we will give an example where
this map is not injective for $n = 1$. Set $S = \Spec(A)$ with
$$
A = k[u, v, s_1, s_2, t_1, t_2]/I
$$
where $k$ is a field and
$$
I = (-us_1 + vt_1 + ut_2, vs_1 + us_2 - vt_2, vs_2, ut_1).
$$
Denote $\overline{u}$ the class of $u$ in $A$ and similarly for
the other variables.
Let $M = (Ax \oplus Ay)/A(\overline{u}x + \overline{v}y)$ so that
$$
\text{Sym}(M) = A[x, y]/(\overline{u}x + \overline{v}y)
= k[x, y, u, v, s_1, s_2, t_1, t_2]/J
$$
where
$$
J = (-us_1 + vt_1 + ut_2, vs_1 + us_2 - vt_2, vs_2, ut_1, ux + vy).
$$
In this case the projective bundle associated to the quasi-coherent
sheaf $\mathcal{E} = \widetilde{M}$ on $S = \Spec(A)$ is the scheme
$$
P =
\text{Proj}(\text{Sym}(M)).
$$
Note that this scheme as an affine open covering
$P = D_{+}(x) \cup D_{+}(y)$.
Consider the element
$m \in M$ which is the image of the element
$us_1x + vt_2y$. Note that
$$
x(us_1x + vt_2y) = (s_1x + s_2y)(ux + vy) \bmod I
$$
and
$$
y(us_1x + vt_2y) = (t_1x + t_2y)(ux + vy) \bmod I.
$$
The first equation implies that $m$ maps to zero as a
section of $\mathcal{O}_P(1)$ on $D_{+}(x)$ and the second
that it maps to zero as a section of $\mathcal{O}_P(1)$ on $D_{+}(y)$.
This shows that $m$ maps to zero in $\Gamma(P, \mathcal{O}_P(1))$.
On the other hand we claim that $m \not = 0$, so that $m$ gives
an example of a nonzero global section of $\mathcal{E}$ mapping to zero
in $\Gamma(P, \mathcal{O}_P(1))$. Assume $m = 0$
to get a contradiction. In this case there exists
an element $f \in k[u, v, s_1, s_2, t_1, t_2]$ such that
$$
us_1x + vt_2y = f(ux + vy) \bmod I
$$
Since $I$ is generated by homogeneous polynomials of degree $2$ we
may decompose $f$ into its homogeneous components and take the
degree 1 component. In other words we may assume that
$$
f = au + bv + \alpha_1s_1 + \alpha_2s_2 + \beta_1t_1 + \beta_2t_2
$$
for some $a, b, \alpha_1, \alpha_2, \beta_1, \beta_2 \in k$.
The resulting conditions are that
$$
\begin{matrix}
us_1 - u(au + bv + \alpha_1s_1 + \alpha_2s_2 + \beta_1t_1 + \beta_2t_2)
\in I \\
vt_2 - v(au + bv + \alpha_1s_1 + \alpha_2s_2 + \beta_1t_1 + \beta_2t_2)
\in I
\end{matrix}
$$
There are no terms $u^2, uv, v^2$ in the generators of $I$ and
hence we see $a = b = 0$. Thus we get the relations
$$
\begin{matrix}
us_1 - u(\alpha_1s_1 + \alpha_2s_2 + \beta_1t_1 + \beta_2t_2)
\in I \\
vt_2 - v(\alpha_1s_1 + \alpha_2s_2 + \beta_1t_1 + \beta_2t_2)
\in I
\end{matrix}
$$
We may use the first generator of $I$ to replace any occurrence of
$us_1$ by $vt_1 + ut_2$, the second generator of $I$ to replace any
occurrence of $vs_1$ by $-us_2 + vt_2$, the third generator
to remove occurrences of $vs_2$ and the third to remove occurrences
of $ut_1$. Then we get the relations
$$
\begin{matrix}
(1 - \alpha_1)vt_1 + (1 - \alpha_1)ut_2 - \alpha_2us_2 - \beta_2ut_2 = 0 \\
(1 - \alpha_1)vt_2 + \alpha_1us_2 - \beta_1vt_1 - \beta_2vt_2 = 0
\end{matrix}
$$
This implies that $\alpha_1$ should be both $0$ and $1$ which is
a contradiction as desired.
\end{example}


\begin{lemma}
\label{lemma-projective-bundle-separated}
Let $S$ be a scheme.
The structure morphism $\mathbf{P}(\mathcal{E}) \to S$ of a
projective bundle over $S$ is separated.
\end{lemma}

\begin{proof}
Immediate from Lemma \ref{lemma-relative-proj-separated}.
\end{proof}

\begin{lemma}
\label{lemma-projective-space-bundle}
Let $S$ be a scheme. Let $n \geq 0$. Then
$\mathbf{P}^n_S$ is a projective bundle over $S$.
\end{lemma}

\begin{proof}
Note that
$$
\mathbf{P}^n_{\mathbf{Z}} =
\text{Proj}(\mathbf{Z}[T_0, \ldots, T_n]) =
\underline{\text{Proj}}_{\Spec(\mathbf{Z})}
\left(\widetilde{\mathbf{Z}[T_0, \ldots, T_n]}\right)
$$
where the grading on the ring $\mathbf{Z}[T_0, \ldots, T_n]$ is given by
$\deg(T_i) = 1$ and the elements of $\mathbf{Z}$ are in degree $0$.
Recall that $\mathbf{P}^n_S$ is defined as
$\mathbf{P}^n_{\mathbf{Z}} \times_{\Spec(\mathbf{Z})} S$.
Moreover, forming the relative homogeneous spectrum commutes with base change,
see Lemma \ref{lemma-relative-proj-base-change}.
For any scheme $g : S \to \Spec(\mathbf{Z})$ we have
$g^*\mathcal{O}_{\Spec(\mathbf{Z})}[T_0, \ldots, T_n]
= \mathcal{O}_S[T_0, \ldots, T_n]$.
Combining the above we see that
$$
\mathbf{P}^n_S = \underline{\text{Proj}}_S(\mathcal{O}_S[T_0, \ldots, T_n]).
$$
Finally, note that
$\mathcal{O}_S[T_0, \ldots, T_n] = \text{Sym}(\mathcal{O}_S^{\oplus n + 1})$.
Hence we see that $\mathbf{P}^n_S$ is a projective bundle over $S$.
\end{proof}



